\section{Uji Validasi Eksternal \textit{Smart Building IoT} (H1)}

\begin{longtable}{
    @{}
    >{\raggedright\arraybackslash}p{3.2cm}          % Pembanding
    >{\raggedleft\arraybackslash}p{2.5cm}           % ΔF1
    >{\raggedleft\arraybackslash}p{3.0cm}           % p (Wilcoxon)
    >{\raggedleft\arraybackslash}p{2.5cm}           % r (rank-biserial)
    @{}
}
\caption{Post-hoc CoDA-EDA vs baseline pada Smart Building IoT  (Friedman $\chi^2 = 21,4$; $p = 0,002$; $n = 5$ seed)} \\
\toprule
Pembanding & $\Delta$F1 & $p$ (Wilcoxon) & $r$ (rank-biserial) \\
\midrule
\endfirsthead
\caption{Post-hoc CoDA-EDA vs baseline pada Smart Building IoT (lanjutan)} \\
Pembanding & $\Delta$F1 & $p$ (Wilcoxon) & $r$ (rank-biserial) \\
\midrule
\endhead
\bottomrule
\endfoot

compact-GA  & $-0,023$ & $0,438$ & $-0,60$ \\
compact-PSO & $-0,026$ & $0,625$ & $-0,60$ \\
iBOA        & $-0,005$ & $1,000$ & $-0,20$ \\
BOA         & $-0,057$ & $0,062$ & $-1,00$ \\
ACO         & $-0,056$ & $0,062$ & $-1,00$ \\
UBK         & $+0,104$ & $0,062$ & $+1,00$ \\

\end{longtable}

\textbf{Catatan:} Dengan hanya lima seed, nilai $p$ terkecil yang mungkin pada uji Wilcoxon dua sisi adalah $2/2^5 = 0,0625$, sehingga taraf $0,05$ tidak dapat dicapai berapa pun besar perbedaannya. 
Hasil eksternal karena itu ditafsirkan secara deskriptif (arah $r$ dan posisi relatif), bukan sebagai vonis signifikansi; uji Friedman omnibus tetap signifikan ($p = 0,002$).